\documentclass[../main.tex]{subfiles}

\makeatletter
\@ifundefined{fromRoot}{%
  \newcommand{\fromRoot}[1]{../#1}
  
  \usepackage{xr}
  \externaldocument{../main}
}{}

\def\input@path{{\subfix{../}}}
%or: \def\input@path{{/path/to/folder/}{/path/to/other/folder/}}
\makeatother

\graphicspath{
  {\subfix{../}}
  {\subfix{./figures}}
  {\subfix{../figures}}
  {\subfix{./figures/logos-thesis/}}
  {\subfix{../figures/logos-thesis/}}
  {\subfix{./figures/rtexps-pics/}}
  {\subfix{../figures/rtexps-pics/}}
}

\hypersetup{
    pdfauthor   = {Camille MONIÈRE},
    pdftitle    = {Th\`{e}se (Présentation: contexte)},
    pdfsubject  = {Th\`{e}se (Présentation: contexte)},
%    pdfkeywords = {mots-cl\'{e}s},
}

\begin{document}

\section{Introduction}

\subsection{Contexte}

\begin{frame}{Télécommunications numériques sans-fil}
  \begin{columns}
    \begin{column}{.4\linewidth}
      \begin{center}
        \placeholder{width=.85\linewidth, height=.6\textheight}
        \captionof{figure}{Illustration : une antenne, deux devices.}
      \end{center}
    \end{column}
    \begin{column}{.55\linewidth}
      \small
      \begin{overlayarea}{\linewidth}{.6\textheight}
        \begin{ctrlitemize}{1em}
          \item Plus fiables et résilients que leurs pendants analogiques.
          \item<2-> Champ regroupant :
          \begin{ctrlitemize}{1ex}
            \scriptsize
            \item<2-> les réseaux cellulaires \\\hspace{1em} \textit{3G, 4G/LTE, 5G},
            \item<3-> les réseaux en champs proches \\\hspace{1em} \textit{ZigBee, NFC, Bluetooth},
            \only<4>{\item et les réseaux longues portées basses puissances \acrshortpl{lpwan} \\\hspace{1em} \textit{LoRa/LoRaWAN, SigFox, NB-IoT}.}
            \only<5->{\item et \textbf{\textcolor{red}{les réseaux longues portées basses puissances \acrshortpl{lpwan}}} \\\hspace{1em} \textit{LoRa/LoRaWAN, SigFox, NB-IoT}.}
          \end{ctrlitemize}
        \end{ctrlitemize}
      \end{overlayarea}
    \end{column}
  \end{columns}
\end{frame}

\begin{frame}{\acrfullpl{lpwan}}
  \small
  \begin{overlayarea}{\linewidth}{.1\textheight}
    \only<1>{
      Fort développement de l'\acrfull{iot} \textit{[internet des objets]} durant la dernière décennie
      \textcolor{red}{$\Rightarrow$} développement des \acrshortpl{lpwan}\cite{IEEEStandardLPWAN}.
    }
    \only<2>{
      Le \acrshortpl{lpwan} ont connu un fort développement durant la dernière décennie, de nombreux standards ont vu le jour\cite{goursaudDedicatedNetworksIoT2015, sinhaSurveyLPWATechnology2017}.%
    }
  \end{overlayarea}

  \begin{overlayarea}{\linewidth}{.65\textheight}
    \only<1>{
      \begin{columns}
        \begin{column}{.4\linewidth}
          \scriptsize
          Caractéristiques des \acrshortpl{lpwan} :
          \begin{ctrlitemize}{4pt}
            \item portées longues,
            \item coûts faibles,
            \item consommations énergétiques réduites.
          \end{ctrlitemize}
          $\Rightarrow$ Permettre le déploiement massif de d'objets connectés alimentés par batterie.\\
          $\Rightarrow$ Requiert des débits plus faibles, que les réseaux cellulaires notamment.
        \end{column}
        \begin{column}{.5\linewidth} \centering
          \includegraphics[width=\linewidth]{figures/drawiopdf/lpwan_and_co}
          \captionof{figure}{Aperçu des caractéristiques de trois technologies utilisées dans les \acrshortpl{lpwan}, SigFox, LoRa et NB-IoT.}
        \end{column}
      \end{columns}
    }
    \only<2>{
      \begin{center}
        \ra{1.3}
        % \begin{noindent}
          \scriptsize
          \begin{tabular}[t]{@{}r@{\phantom{XXX}}r@{\phantom{XX}}r@{\phantom{XX}}r@{}}
            \toprule
                                   & \textbf{SigFox}   & \textbf{LoRa}     & \textbf{NB-IoT}     \\ \midrule
            \textbf{Modulation}    & \acrshort{bpsk}   & \acrshort{css}    & \acrshort{qpsk}     \\ 
            \textbf{Fréquence centrale}
                                   & $868$ MHz         & $868$ MHz         & bandes LTE          \\
            \textbf{Largeur de bande}
                                   & $0.1$ kHz         & $125$/$250$ kHz   & $200$ kHz           \\
            \textbf{Débit max de données} 
                                   & $0.1$ kb/s        & $50$ kb/s         & $200$ kb/s          \\
            \textbf{Taille max de la charge utile}
                                   & $96$ bits         & $1944$ bits       & $12.8$ kilobits     \\
            \textbf{Portée urbaine}
                                   & $10$ km           & $5$ km            & $1$ km              \\
            \textbf{Portée rurale} & $40$ km           & $20$ km           & $10$ km             \\
            \textbf{Résilience aux interférences}
                                   & très grande       & très grande       & basse               \\
            \textbf{Code correcteur d'erreurs}
                                   & aucun             & codes de Hamming  & codes modernes      \\
            \bottomrule
          \end{tabular}
          % \end{noindent}
        \captionof{table}{Aperçu des caractéristiques de trois technologies utilisées dans les \acrshortpl{lpwan}, SigFox, LoRa et NB-IoT.}
        \ra{1}
      \end{center}
    }
  \end{overlayarea}
  \only<1>{\blfootnotetext{\textcite{IEEEStandardLPWAN}}}
  \only<2>{
    \blfootnotetext{\textcite{goursaudDedicatedNetworksIoT2015, sinhaSurveyLPWATechnology2017}}
  }
\end{frame}


\begin{frame}{Problème des préambules}
  % \begin{overlayarea}{\linewidth}{.05\textheight}
  % \end{overlayarea}
  \begin{columns}
    \begin{column}{.25\linewidth}
      \begin{overlayarea}{\linewidth}{.45\textheight}
        \includegraphics[width=\linewidth]{figures/tikzpicture/osi_layers_stdl}
      \end{overlayarea}
    \end{column}
    \begin{column}{.02\linewidth}
    \end{column}
    \begin{column}{.69\linewidth}
      \begin{overlayarea}{\linewidth}{.45\textheight}
        \only<1,2>{
          \includegraphics[width=\linewidth]{figures/tikzpicture/joint_frame_tikz_bare.pdf}
        }
        \only<3>{
          \includegraphics[width=\linewidth]{figures/tikzpicture/joint_frame_tikz.pdf}
        }
      \end{overlayarea}
    \end{column}
  \end{columns}

  \begin{overlayarea}{\linewidth}{.4\textheight}
    \only<1>{
      \includegraphics[width=\linewidth]{figures/tikzpicture/joint_frame_tikz_leg.pdf}
    }
    \only<2->{ \centering
    \includegraphics[width=\linewidth]{figures/tikzpicture/joint_frame_tikz_leg-2.pdf}\\
    {\tiny \fullcite{polyanskiyAsynchronousCommunicationExact2013}}
    }
  \end{overlayarea}
\end{frame}


\subsection[Le projet \acrshort{qcsp}]{Le projet \acrfull{qcsp}}

\begin{frame}{\subsecname}{\href{https://qcsp.univ-ubs.fr/}{https://qcsp.univ-ubs.fr/}}

  \begin{columns}
    \begin{column}{.4\linewidth} \centering
      \begin{columns}
        \begin{column}{.4\linewidth}
          \includegraphics[width=\textwidth]{QCSP_Logo.pdf}
        \end{column}
        \begin{column}{.2\linewidth}
          \includegraphics[width=\linewidth]{anr_approved.jpg}
        \end{column}
      \end{columns}
      \small \vspace{5pt} Projet ANR \texttt{ANR-19-CE25-0013-01}
    \end{column}
    \begin{column}{.5\linewidth}
      \emph{``L'objectif du projet QCSP est de contribuer à l'évolution des réseaux IoT en établissant,
        en mettant en œuvre et testant un nouveau type de modulation codée.''}
    \end{column}
  \end{columns}

  \vspace{2 em}

  \begin{center}
    Partenaires :\\
    \includegraphics[width=0.6\linewidth]{figures/logos-thesis/partners-logos.png}
  \end{center}
\end{frame}

\begin{frame}{Idées de \acrshort{qcsp}}
  \begin{overlayarea}{\linewidth}{0.78\textheight}
    \small
    \begin{columns}
      \begin{column}<1,3>{.4\linewidth} \centering
        \textsc{Internet de Objets (IoT)}
      \end{column}
      \begin{column}{.05\linewidth} \centering
        \only<1,2>{\phantom{VS}}
        \only<3>{VS}
      \end{column}
      \begin{column}<2,3>{.4\linewidth} \centering
        \textsc{Les technologies spatiales}
      \end{column}
    \end{columns}

    \begin{columns}
      \begin{column}<1,3>[t]{.45\linewidth}
        \vspace{-2 ex}
        \begin{ctrlitemize}{-5pt}
          \item Croissance exponentielle ces 10 dernières années,
          \item 50 milliards d'objets connectés attendus,
        \end{ctrlitemize}

        \vspace{-1 ex}
        \hspace{7 em} \emph{et pourtant\dots}
        \vspace{-1 ex}

        \begin{ctrlitemize}{-2pt}
          \item les métadonnées de détection et synchronisation représentent parfois jusqu'à  50\% de la consommation de ressources \cite{durisiMassiveUltrareliableLowLatency2016}.
        \end{ctrlitemize}
        \begin{center}
          \vspace{-1 em}
          \includegraphics<1,3>[width=.4\linewidth, clip]{Anonymous_Factory.png}
        \end{center}
      \end{column}
      \begin{column}{.05\linewidth}
        \hfill
      \end{column}
      \begin{column}<2,3>[t]{.45\linewidth}
        \vspace{-1 ex}
        \begin{ctrlitemize}{-5pt}
          \item Privilégie la stabilité
        \end{ctrlitemize}

        \vspace{-1 ex} \hspace{7 em} \emph{et\dots}
        \vspace{-1 ex}

        \begin{ctrlitemize}{-5pt}
          \item le Cyclic-Code Shift Keying (CCSK) est utilisé dans \textit{Quasi-Zenith Satellite System}\\ % dans la bande LEX
          {\tiny  (système japonais d'amélioration du ``\emph{GPS}'')} \\
          \item des codes correcteurs non-binaires sont utilisé dans BeiDou \cite{openservicesignalb2bBeiDouNavigationSatellite2020}\\ {\tiny  (Le ``\emph{GPS}'' chinois.)}
        \end{ctrlitemize}

        \begin{center}
          \includegraphics<2,3>[width=.4\linewidth]{pngegg.png}
        \end{center}
      \end{column}
    \end{columns}
  \end{overlayarea}
  \only<1>{\blfootnotetext{\textcite{durisiMassiveUltrareliableLowLatency2016}}}
  \only<2>{\blfootnotetext{\textcite{openservicesignalb2bBeiDouNavigationSatellite2020}}}
  \only<3>{\blfootnotetext{\textcite{durisiMassiveUltrareliableLowLatency2016, openservicesignalb2bBeiDouNavigationSatellite2020}}}
\end{frame}


\begin{frame}{\acrfull{ccsk}}
  {}
  \begin{overlayarea}{\linewidth}{.02 \textheight}
  \end{overlayarea}
  \begin{overlayarea}{\linewidth}{.18 \textheight}
    \centering
    \normalsize Modulation d'un mot de code de symboles de $p$ bits par la rotation circulaire d'une séquence de pseudo-bruit binaire (\acrshort{pn}) \pn{0}\cite{dillardCyclicCodeShift2003} de $2^p = q$ chips \\\emph{\textsf{bijection de \gf{q}\footnotemark vers l'ensemble des rotations circulaires de \pn{0}}}.
    \vfill
  \end{overlayarea}

  \begin{overlayarea}{\linewidth}{.4 \textheight}
    \begin{columns}
      \begin{column}{.3\linewidth} \centering
        \centering

        \begin{tabular}{@{}r r l@{}}
          \toprule
          \multicolumn{3}{c}{\textbf{Paramètres}} \\ \midrule
          $p$    & $=$ & $2$                      \\
          $q$    & $=$ & $4$                      \\
          \pn{0} & $=$ & \Ob{}\Xb{}\Xb{}\Xb{}     \\
          \bottomrule
        \end{tabular}

        \Ob{} $= +1$, \Xb{} $= -1$

      \end{column}
      \begin{column}{.3\linewidth}
        \begin{tabular}{@{}r c l@{}}
          \toprule
          Symbole & Notation & Rotation             \\ \midrule
          $0$     & \pn{0}   & \Ob{}\Xb{}\Xb{}\Xb{} \\
          $1$     & \pn{1}   & \Xb{}\Ob{}\Xb{}\Xb{} \\
          $2$     & \pn{2}   & \Xb{}\Xb{}\Ob{}\Xb{} \\
          $3$     & \pn{3}   & \Xb{}\Xb{}\Xb{}\Ob{} \\
          \bottomrule
        \end{tabular}
      \end{column}
    \end{columns}
  \end{overlayarea}


  \begin{overlayarea}{\linewidth}{.2 \textheight}
    \begin{columns}
      \begin{column}{.3\linewidth} \centering
        Mot de code $\vect{C}$\\$\{ 0, \; 3, \; 1,\; 2\}$
      \end{column}
      \begin{column}{.1\linewidth} \centering
        $\Rightarrow$
      \end{column}
      \begin{column}{.55\linewidth} \centering
        Trame \acrshort{ccsk} $\vect{F}_{CCSK}$\\$\{ \mpn{0} | \mpn{3} | \mpn{1} | \mpn{2} \}$ = \Ob{}\Xb{}\Xb{}\Xb{}~~\Xb{}\Xb{}\Xb{}\Ob{}~~\Xb{}\Ob{}\Xb{}\Xb{}~~\Xb{}\Xb{}\Ob{}\Xb{}
      \end{column}
    \end{columns}

  \end{overlayarea}
  \blfootnotetext{\textcite{dillardCyclicCodeShift2003}}
  \footnotetext{Corps de Galois d'ordre $p$}
\end{frame}


\begin{frame}{\acrshort{ccsk} : Démodulation}
  
\end{frame}

\begin{frame}{\acrfull{ldpc}}
  \textbf{TODO, Juste le principe}
\end{frame}


\end{document}
